%% Presentation speech for the OrdinalFed talk (IEEE WIECON-ECE 2026, Paper 518) -- simple version
\documentclass[12pt,a4paper]{article}
\usepackage[margin=2.2cm]{geometry}
\usepackage[T1]{fontenc}
\usepackage[scaled]{helvet}
\renewcommand{\familydefault}{\sfdefault}
\usepackage{xcolor}
\usepackage{enumitem}
\usepackage[hidelinks,pdftitle={OrdinalFed - Presentation Speech}]{hyperref}
\definecolor{navy}{HTML}{14213D}
\definecolor{accent}{HTML}{2F6FB3}
\definecolor{muted}{HTML}{5B6575}
\setlength{\parindent}{0pt}
\setlength{\parskip}{6pt}
\linespread{1.3}
\pagestyle{plain}

% \slide{number}{title}{time}
\newcommand{\slide}[3]{%
  \par\bigskip\noindent{\color{accent}\large\bfseries Slide #1 \textperiodcentered\ #2}%
  \hfill{\color{muted}\small #3}\par\smallskip}
% stage direction
\newcommand{\act}[1]{{\color{muted}\itshape (#1)}}

\begin{document}

{\color{navy}\LARGE\bfseries OrdinalFed: Presentation Speech}\par\smallskip
{\color{muted}IEEE WIECON-ECE 2026 \textperiodcentered\ Paper ID 518 \textperiodcentered\ Presenter: Md Tamim Hossain}

\medskip
\fcolorbox{accent}{accent!6}{\parbox{\dimexpr\linewidth-2\fboxsep-2\fboxrule}{%
\textbf{Total time:} about 9--10 minutes for 17 slides.\\
\textbf{How to speak:} slowly, in short sentences. Look at the audience. Pause after each number.
The grey text in brackets tells you what to do, not what to say.}}

\slide{1}{Title}{0:20}
Good morning, everyone. My name is Md Tamim Hossain, from the CSE department of MIST.
Today I will present our paper, OrdinalFed. It is about training an eye-disease model across
many hospitals. This is joint work with Najifa Raksanda, M.\,M. Tamim Sharif and Aria
Nawshin Ashka.

\slide{2}{Diabetic retinopathy}{0:40}
Diabetes can damage the blood vessels in the eye. This is called diabetic retinopathy, or DR.
It is a major cause of blindness that we can prevent. Bangladesh has more than 13 million
people with diabetes. But there are very few eye doctors outside the big cities.

Doctors grade DR from 0, which is healthy, to 4, which is very severe. If we find Grade 2 or
higher early, we can save a patient's sight. So a computer that grades DR automatically would
help a lot.

\slide{3}{Why federated learning?}{0:40}
To train a good model, we need images from many hospitals. But hospitals cannot share patient
images, because of privacy. Federated learning solves this.

Each hospital trains the model on its own images. Then it sends only the model, not the images,
to a central server. The server combines all the models into one and sends it back. This repeats
for several rounds. The usual way to combine them is called FedAvg. In FedAvg, a hospital with
more images gets a bigger vote.

\slide{4}{Two problems}{0:45}
But there are two problems. First, hospitals see very different patients. Some see mostly
healthy eyes. Others see mostly serious cases. When the data are this uneven, FedAvg can break
down. In our tests, its score fell from 0.647 to 0.378.

\act{Point to the small note at the bottom.} We use a score called QWK. It checks how close the
predicted grade is to the true grade. One means perfect.

Second, the usual training method treats every mistake the same. Calling a Grade 4 eye
``Grade 3'' is a small mistake. Calling it ``Grade 0'' is a big one. But the usual method sees
no difference.

\slide{5}{Example: three very different hospitals}{0:50}
Here is a simple example. Think of three hospitals. A big city hospital has many patients, but
most have healthy eyes. An eye hospital has many serious cases. A village clinic has only a few
patients.

FedAvg gives votes by the number of images. So the city hospital gets almost all the votes. The
model learns that most eyes are healthy. Then it starts missing the serious cases and the early
cases.

These bars are just a picture of the idea. But we measured the same thing: the two smallest
hospitals got only 5.9 percent of the vote.

\slide{6}{OrdinalFed has two parts}{0:45}
Our method, OrdinalFed, has two parts.

Part one is at each hospital. We add an extra penalty when the predicted grade is far from the
true grade. So a far-off mistake costs much more than a near one.

Part two is at the server. The server keeps 50 test images, 10 of each grade. After each round,
it tests every hospital's model on these images. Better models get a bigger vote. The hospitals
do not need to change anything.

\slide{7}{How the server sets the weights}{0:35}
This is how the server sets the votes. First, it gives each model a score on its 50 images.
Second, it turns the scores into votes. Third, it mixes these votes with the image counts.

At the start, both count equally. Over the rounds, the image counts matter less and less. At the
end, only the scores matter. No hospital is ever left out.

\slide{8}{Same exam for every hospital}{0:40}
So, back to our three hospitals. You can think of it like a school exam. Every hospital takes
the same exam. The vote depends on how well you did, not just on how big you are. And every
hospital keeps a small minimum vote.

In our tests, the vote of the two smallest hospitals went up from 5.9 percent to 40.8 percent.
And the worst score went up from 0.378 to 0.514. The hospitals did not change anything.

\slide{9}{Experimental setup}{0:30}
Now, our experiments. We used the DDR dataset, with more than 12,000 eye images. We had five
hospitals and used the EfficientNet-B3 model. We trained for 10 rounds and repeated every test
three times. We compared our method with FedAvg, FedProx and two other versions. We also tested
all models on a second dataset, APTOS 2019.

\slide{10}{Three ways to split the data}{0:35}
We split the images among the five hospitals in three ways. Each bar is one hospital. The
colours show which grades it has.

\act{Point left.} On the left, every hospital gets a similar mix. That is the easy case.
\act{Point to the middle.} In the middle, each hospital gets mostly one grade. That is the hard
case. One hospital has only 32 images, all healthy. \act{Point right.} On the right, the split
copies five real kinds of clinics.

\slide{11}{Test QWK after every round}{0:35}
This chart shows the score after each round. Orange is FedAvg. Green is FedProx. Blue is our
method.

\act{Point to the middle panel.} Look at the middle, the hard case. FedAvg crashes after round 5.
Our method does not crash, but it goes up and down. FedProx stays steady. On the left and right,
all methods do about the same.

\slide{12}{Worst round under extreme skew}{0:30}
Here is the worst score of each method in the hard case. Higher is better. FedAvg falls to
0.378. Our method keeps it at 0.514. FedProx is the most steady, at 0.641. Our method plus
FedProx gets close, at 0.627.

\slide{13}{Why does it work?}{0:45}
Why does our method help? We tested each part on its own.

The new penalty alone did not help. Its score, 0.380, was the same as FedAvg. Votes based on
scores helped: 0.586. But giving every hospital an equal vote helped even more: 0.630.

So the real reason is simple. Small hospitals get a bigger voice. With only 50 test images, the
scores are too noisy to rank the hospitals well.

\slide{14}{OrdinalFed or FedProx?}{0:50}
So why use our method instead of FedProx? The difference is where the change happens. FedProx
changes the training at every hospital. Each hospital must change its code and tune a setting.
Our method changes only the server. The hospitals do nothing new.

To be honest, FedProx was more steady in our tests. But our method is useful in three cases.
One: when you can only change the server. Two: when hospital sizes are very different. Three:
when finding rare, serious cases matters. Our method found more Grade 3 cases: 0.36 against
0.31. And you can use both together. That reached 0.627.

\slide{15}{More results}{0:30}
A few more results. When the data are not very uneven, all methods score about the same. Our
method found the most Grade 3 cases. On the second dataset, APTOS, all models still worked well,
close to a model trained on all the data. And our method adds only about one second per round at
the server.

\slide{16}{Takeaways}{0:40}
To sum up. One: our method stops most of FedAvg's crash, and the hospitals change nothing. Two:
it works because small hospitals get a bigger voice. Three: FedProx is still the most steady
method.

Our work has limits. The 50 test images are few. Grade 1 is still hard to detect. And we used one
dataset with small images. Next, we want to try more test images, bigger images and more
datasets.

\slide{17}{Thank you}{0:05}
Thank you for listening. I am happy to answer your questions.

\bigskip
\noindent{\color{navy}\large\bfseries If you are short of time}\par
Make Slides 9, 10 and 15 very short (one or two sentences each). Never skip Slides 5, 8 and
11--14: they tell the main story.

\bigskip
\noindent{\color{navy}\large\bfseries Likely questions and simple answers}
\begin{enumerate}[leftmargin=*,itemsep=6pt]
  \item \textbf{Why not just use FedProx?}\\
  FedProx was more steady in our tests. But our method changes only the server, so it helps when
  you cannot change every hospital's software. And the two can be used together: that reached
  0.627. (Slide 14.)
  \item \textbf{Is keeping 50 labelled images on the server a privacy problem?}\\
  In our tests, these images came from a public dataset, not from the hospitals. In real use,
  the server needs a few labelled images. This is one of our limits.
  \item \textbf{Why do equal votes beat score-based votes?}\\
  In the hard case, FedAvg lets a few big hospitals take over. Any rule that gives small
  hospitals more say helps. With only 50 images, the scores are too noisy.
  \item \textbf{Why is the test score (0.682) lower than the validation score (0.834)?}\\
  The test images came from different cameras, and many Grade 2 eyes were called Grade 0. This
  affected all methods in the same way.
  \item \textbf{Why is Grade 1 so hard?}\\
  Grade 1 is only 5\% of our training images, and its signs are very small at 224 $\times$ 224
  pixels. Bigger images may help.
  \item \textbf{So what does the extra penalty add?}\\
  In normal (non-federated) training it helps. In our federated tests, it did not add much. We
  report this honestly.
\end{enumerate}

\end{document}
